\section{三类典型方程行为迥异}
\label{sec:15-Differential-Equations-and-Dynamical-Systems/03-PDE-Introduction/02}

把同一个带尖角的初始起伏分别交给三个方程：热方程 \(\partial_tu=\kappa\partial_{xx}u\)、波动方程 \(\partial_{tt}u=c^2\partial_{xx}u\)、Laplace 方程 \(\Delta u=0\)。三者的空间项写出来几乎一样，区别只在时间项——一阶、二阶、干脆没有。就这点差别，三个解的下场完全不同：一个把尖角磨没，一个带着尖角原样跑，第三个连``下场''都谈不上，因为它没有时间可走。本节要做的，是把这三种形态描成一眼能认出的样子。

\begin{center}
\begin{tikzpicture}[>=stealth]
  \begin{scope}
    \draw[->] (0,0) -- (3.2,0) node[right] {\small \(x\)};
    \draw[thick] plot[smooth] coordinates
      {(0,0.10) (0.4,0.12) (0.6,1.65) (0.8,1.15) (1.0,1.80)
       (1.2,1.25) (1.4,1.75) (1.6,1.20) (1.8,1.70) (2.0,0.14)
       (2.4,0.10) (3.0,0.10)};
    \draw[dashed] plot[smooth] coordinates
      {(0,0.16) (0.5,0.34) (0.9,1.12) (1.3,1.32) (1.7,1.20)
       (2.1,0.52) (2.6,0.22) (3.0,0.18)};
    \draw[dotted,thick] plot[smooth] coordinates
      {(0,0.28) (0.75,0.62) (1.5,0.80) (2.25,0.60) (3.0,0.30)};
    \node at (1.5,2.50) {\small 抛物型：抹平};
  \end{scope}

  \begin{scope}[shift={(4.4,0)}]
    \draw[->] (0,0) -- (3.2,0) node[right] {\small \(x\)};
    \draw[thick] (0.25,0) -- (0.6,1.4) -- (0.95,0);
    \draw[dashed] (1.15,0) -- (1.5,1.4) -- (1.85,0);
    \draw[dotted,thick] (2.05,0) -- (2.4,1.4) -- (2.75,0);
    \draw[->] (0.75,1.65) -- (1.75,1.65);
    \node at (1.25,1.9) {\small 速度 \(c\)};
    \node at (1.5,2.50) {\small 双曲型：搬运};
  \end{scope}

  \begin{scope}[shift={(8.8,0)}]
    \draw[thick] (0.2,0.1) rectangle (2.8,1.8);
    \foreach \s in {0.75,1.25,1.75,2.25}{
      \draw[gray] plot[smooth] coordinates
        {(\s,0.1) ({\s+0.18},0.95) (\s,1.8)};
    }
    \node at (0.05,0.95) {\small \(0\)};
    \node at (3.0,0.95) {\small \(1\)};
    \node at (1.5,2.50) {\small 椭圆型：平衡};
  \end{scope}
\end{tikzpicture}
\end{center}

\emph{同一类起伏交给三种方程：抹平、原样搬运、以及由边界一次锁死的静态平衡。}

\subsection{热方程把高频按平方速率处死}

铁棒一端 100 度、另一端 0 度，中间那道近乎垂直的跳变，几秒后就模糊了；再等一会儿，整根棒趋向一条温和的斜线。这不是物理层面的比喻，方程自己写着答案。在周期边界下试一个单频分量 \(u(t,x)=T(t)\sin(kx)\)，代入 \(\partial_tu=\kappa\partial_{xx}u\)：

\[
\begin{aligned}
T'(t)\sin(kx)&=-\kappa k^2T(t)\sin(kx),\\
T(t)&=e^{-\kappa k^2t}\,T(0).
\end{aligned}
\]

衰减率是 \(\kappa k^2\)，随波数平方增长。\(k=10\) 的分量衰减得比 \(k=1\) 快一百倍，\(k=100\) 快一万倍。图中第一格那条毛糙的实线之所以在下一个时刻就变成平顺的虚线，原因全在这个平方上：细节属于大 \(k\)，大 \(k\) 死得最快。

由此得到一条颇为极端的结论。\hyperref[sec:08-Numerical-Analysis/08-Fourier-and-Spectral-Methods/01]{``从 Fourier 级数到连续变换''}里，光滑的函数高频系数掉得快，那是被动的对应关系；热方程则是主动施加这件事——不管初值有多粗糙，只要 \(t>0\)，所有高频都已经被指数压掉，解立刻是 \(C^\infty\) 的。一个分段常值的初始温度，过一瞬间就无限可微。

代价写在同一个公式里：\(e^{-\kappa k^2t}\) 把高频压到远低于机器精度之后，那部分信息就再也拿不回来了。热方程的时间箭头是单向的。这一点在扩散模型的前向过程里被直接当作工具使用——\hyperref[sec:06-Random-Processes/06-Diffusion-and-SDE/02]{``Fokker--Planck 方程描述密度的流动''}给出的密度演化方程恰是带漂移的热方程，前向加噪之所以能把任何数据分布磨成高斯，靠的正是这种不可逆的抹平；而所谓逆向生成并不是把热方程倒着解，是另学一个模型把丢掉的东西补回来。

最大值原理从另一个角度锁住解的形状：无热源时，杆上任一时刻的最高温度只能来自初始分布或者两端边界，内部不会凭空冒出一个比四周都热的点。这条结论有严格证明，也是唯一性与稳定性论证反复使用的工具。

\subsection{波动方程不生产也不摧毁光滑性}

拨一下吉他弦。指尖捏出的三角形带一个尖角，松手后它分成两半各奔一端，撞到琴桥反弹回来。整个过程里那个尖角一直是尖的——没变圆，没变模糊，也没变矮。

d'Alembert 公式把这件事写死了：

\[
u(t,x)=F(x-ct)+G(x+ct).
\]

代进方程验证即可，\(\partial_{tt}u=c^2(F''+G'')\) 与 \(c^2\partial_{xx}u\) 逐项相等。解是两个形状的平移，一个向右一个向左，速度恒为 \(c\)，形状本身不参与演化。图中第二格三个全等的三角形画的就是这件事：横坐标在变，轮廓不变。

与热方程的对照有两处值得记住。其一是传播速度有限：\(x=0\) 处的扰动要等 \(t=L/c\) 才能影响到距离 \(L\) 之外，特征线 \(x\pm ct=\)\,常数之外的区域尚不知道发生了什么。热方程没有这条限制，初始扰动在任意 \(t>0\) 对全空间都有非零影响，尽管远处小到荒谬。其二是正则性守恒：初值有尖角，解就一路带着尖角；初值光滑，解就一路光滑。热方程制造光滑性，波动方程只负责搬运。

往池塘里丢一块石头是同一件事的二维版：波前清晰地向外推，波前之外的水面还是平的。

\subsection{Laplace 方程里内部没有自留地}

\(\Delta u=0\) 没有时间变量，它描述的是已经停下来的东西：稳态温度场、静电势、不可压无旋流的速度势。两条性质合起来说清了``平衡''的含义。

一条是均值性质。调和函数在任一点的值等于以该点为心的任意球面上的平均值。一维时它退化成最朴素的形式：\(\partial_{xx}u=0\) 的解是直线，中点的值就是两端的平均。由此推出内部不能有严格的局部极值——若某点比周围平均更热，热量就会往外流，那就不是稳态了。极值只能落在边界上。

另一条是边界唯一决定内部。给定边界上的值，调和函数在整个区域内部被完全确定，一点多余的自由度都没有。绷在弯折铁框上的肥皂膜是现成的例子：膜的形状你无法从中间去捏，能动的只有边框。图中第三格那些从左壁弯向右壁的灰线就是等值线，它们的走向由左右两侧的赋值一次性定死。

这两条与热方程的最大值原理同根。热方程多出一个时间方向，于是``边界''被拆成两截：初始条件是时间的边界，空间边界条件是空间的边界。

\subsection{判别式一个符号，定下解的全部脾气}

对二阶线性方程 \(au_{xx}+2bu_{xy}+cu_{yy}+\cdots=0\)，判别式 \(b^2-ac\) 的符号给出分类：负是椭圆型，零是抛物型，正是双曲型。分类不是给方程贴标签用的，它一口气回答了四个实际问题——该配什么定解数据，解会有多光滑，扰动怎么传播，以及哪个方向上问题会变坏。

椭圆型没有时间方向，所有空间方向地位平等，该给的是边界条件；给它一个``初始条件''属于类型错误。抛物型时间一阶空间二阶，时间方向是特殊的，向前抹平、向后放大，该给的是初值加边界条件。双曲型时间与空间都是二阶，两者地位接近，扰动沿特征线以有限速度走，该给的是初始位置\emph{和}初始速度——少了后者，弦松手那一刻往哪边走就没人知道。

数值方法的分野同样沿着这条线。抛物方程的显式格式受制于 \(\Delta t\lesssim\Delta x^2/(2\kappa)\)：网格加密一倍，时间步要砍到四分之一，因为离散 Laplacian 的特征值最大到 \(O(\Delta x^{-2})\)，这是一个典型的刚性系统，实践中改用隐式格式或 Crank--Nicolson，理由与 \hyperref[sec:08-Numerical-Analysis/06-ODE/03]{``稳定区域与刚性方程''}中选择 \(A\)-稳定方法的理由一字不差。双曲方程的限制温和得多但更硬性，CFL 条件 \(c\Delta t\le\Delta x\) 要求数值格式的依赖域覆盖住真实的特征依赖域——步子迈得比波还快，格式就在从未被通知的地方取值，必崩。椭圆方程没有时间步可言，离散之后是一次性的大型稀疏线性系统，工作量转移到\hyperref[sec:08-Numerical-Analysis/05-Numerical-Linear-Algebra/04]{``大型稀疏系统与 Krylov 迭代''}那一侧，网格越细矩阵条件数越差，预条件子成为主要战场。

变系数与非线性方程的类型可以随位置改变。跨声速流动里，亚声速区表现为椭圆（扰动向四面八方传），超声速区表现为双曲（扰动只沿 Mach 锥单向传），声速线附近类型退化，单一格式很难全域稳定。流体仿真的多数硬骨头最后都追到这里。

反过来读也成立，而且在调试时更有用。迭代不收敛、残差在全域一起震荡，多半是椭圆问题上用了纯显式的时间推进；解在某个方向上一步算错就再也回不来，多半是双曲系统没有尊重特征的来向；对末端数据的微小扰动敏感到离谱，几乎可以断定是在倒着解抛物方程。数值崩溃很少源于某个加减号写反，更常见的是格式与方程类型没对上。

最后那句``倒着解抛物方程''值得单独算一笔账。下一节把``适定''拆成存在、唯一、连续依赖三条，看看第三条是怎样在热方程的时间反演上崩掉的，以及这个崩法与条件数是同一件事的两种说法。
